% =====================================================================
% pertemuan-04.tex
% Pertemuan 4 - Finite Automata: NFA, DFA, Konversi, dan Minimalisasi
% Sub-CPMK 2.2, 2.3 - Praktik terbimbing
% =====================================================================
\documentclass[aspectratio=169,10pt,t]{beamer}
\usepackage{config/beamerthemeubb}

\UBBmeeting{4}

\begin{document}

\UBBtitleframe{4}{Finite Automata: NFA, DFA, Konversi, dan Minimalisasi}{Praktik terbimbing}

% ---------------------------------------------------------------------
\begin{frame}{Capaian Pertemuan}
  \textbf{CPMK induk:} CPMK 02 --- Memahami konsep Otomata Hingga
  (\textit{Finite Automata}).

  \vspace{3pt}
  \begin{block}{Sub-CPMK yang dinilai pada pertemuan ini}
    \footnotesize
    \begin{itemize}
      \setlength{\itemsep}{2pt}
      \item \textbf{Sub-CPMK 2.2:} Mengimplementasikan simulator NFA dan
            DFA dalam bahasa C untuk mengenali token, dengan berkas kode
            dan keluaran pengujian sebagai bukti.
      \item \textbf{Sub-CPMK 2.3:} Membuat tabel konversi NFA ke DFA dan
            tabel minimalisasi DFA, serta menguji bahwa keduanya
            menerima bahasa token yang sama.
    \end{itemize}
  \end{block}
\end{frame}

% ---------------------------------------------------------------------
\begin{frame}{Agenda Pertemuan (150 menit)}
  \small
  \begin{tabular}{@{}p{1.4cm}p{12.3cm}@{}}
    \toprule
    \textbf{Menit} & \textbf{Kegiatan} \\
    \midrule
    25 & Ceramah: DFA, NFA, dan \textit{epsilon closure} \\
    20 & Demonstrasi: menelusuri string pada DFA dan NFA \\
    55 & Praktik: \textit{subset construction} dan tabel konversi \\
    30 & Praktik: minimalisasi dan tabel partisi \\
    20 & Simulator C, pengujian, dan pemeriksaan bersama \\
    \bottomrule
  \end{tabular}

  \vspace{6pt}
  {\scriptsize\color{ubbblue!75}Porsi praktik 105 menit lebih besar
  daripada ceramah 25 menit. Setiap mahasiswa mengumpulkan tabel dan
  simulator pada akhir pertemuan.}
\end{frame}

% ---------------------------------------------------------------------
\begin{frame}{Posisi Pertemuan pada Rantai Proyek}
  \UBBchain{1}

  \vspace{2pt}
  \begin{itemize}
    \item Pertemuan ini memperkuat simpul \textbf{lexer}: dari pattern
          menuju mesin pengenal.
    \item DFA yang dibangun hari ini menjadi inti lexer tangan pada
          Pertemuan 5.
    \item Tabel konversi dan minimalisasi adalah artefak Sub-CPMK 2.2
          dan 2.3 yang dinilai pada Praktikum 1.
  \end{itemize}
\end{frame}

% ---------------------------------------------------------------------
\begin{frame}{DFA: Definisi}
  \begin{block}{\textit{Deterministic Finite Automaton} (DFA)}
    \small
    \(M = (Q, \Sigma, \delta, q_0, F)\), dengan \(Q\) himpunan state
    berhingga, \(\Sigma\) alfabet, \(\delta : Q \times \Sigma \to Q\)
    fungsi transisi total, \(q_0\) state awal, dan \(F \subseteq Q\)
    himpunan state penerima.
  \end{block}

  \begin{itemize}
    \item Untuk setiap pasangan state dan simbol ada tepat satu state
          tujuan.
    \item Tidak ada transisi \(\varepsilon\) dan tidak ada pilihan
          ganda.
    \item Simulasi hanya menyimpan satu state saat ini.
    \item Kompleksitas waktu \(O(n)\), ruang \(O(1)\) di luar tabel.
  \end{itemize}
\end{frame}

% ---------------------------------------------------------------------
\begin{frame}{Contoh DFA: Jumlah Huruf \texttt{a} Genap}
  \begin{columns}[T,onlytextwidth]
    \begin{column}{0.56\textwidth}
      \centering
      \begin{tikzpicture}[
          state/.style={circle,draw=ubbblue,fill=ubbblue!7,
            minimum size=0.78cm,inner sep=0pt,font=\small},
          accepting/.style={state,double,double distance=1.4pt},
          ar/.style={-{Stealth[length=1.8mm]},draw=ubbblue!70,thick}
        ]
        \node[accepting] (q0) at (0,0) {\(q_0\)};
        \node[state] (q1) at (3.0,0) {\(q_1\)};
        \draw[ar] ([xshift=-0.75cm]q0.west) -- (q0.west);
        \draw[ar] (q0) edge[loop above] node{\scriptsize \(b\)} (q0);
        \draw[ar] (q1) edge[loop above] node{\scriptsize \(b\)} (q1);
        \draw[ar,bend left=18] (q0) to node[above]{\scriptsize \(a\)} (q1);
        \draw[ar,bend left=18] (q1) to node[below]{\scriptsize \(a\)} (q0);
      \end{tikzpicture}
      \par {\scriptsize \(L=\{w \mid\) jumlah \(a\) genap\(\}\)}
    \end{column}
    \begin{column}{0.40\textwidth}
      \small
      \begin{tabular}{@{}ccc@{}}
        \toprule
        \textbf{State} & \textbf{\(a\)} & \textbf{\(b\)} \\
        \midrule
        \(\to *q_0\) & \(q_1\) & \(q_0\) \\
        \(q_1\) & \(q_0\) & \(q_1\) \\
        \bottomrule
      \end{tabular}
      \par\vspace{4pt}
      {\footnotesize Tanda \(\to\) adalah state awal; \(*\) menandai
      state penerima.}
    \end{column}
  \end{columns}
\end{frame}

% ---------------------------------------------------------------------
\begin{frame}{Simulasi DFA}
  \begin{block}{Fungsi transisi diperluas}
    \small
    \(\hat{\delta}(q,\varepsilon)=q\), \quad
    \(\hat{\delta}(q,wa)=\delta(\hat{\delta}(q,w),a)\).
    String \(w\) diterima jika \(\hat{\delta}(q_0,w) \in F\).
  \end{block}

  \begin{enumerate}
    \setlength{\itemsep}{2pt}
    \item Mulai dari \(q_0\).
    \item Untuk setiap simbol masukan, pindah ke \(\delta(q,a)\).
    \item Jika transisi tidak terdefinisi, tolak.
    \item Setelah masukan habis, terima jika state berada di \(F\).
  \end{enumerate}

  \vspace{3pt}
  {\footnotesize Contoh: \(abba\) pada DFA sebelumnya berakhir di
  \(q_0\), sehingga diterima karena jumlah \(a\) genap (dua).}
\end{frame}

% ---------------------------------------------------------------------
\begin{frame}[fragile]{Simulator DFA dalam Bahasa C}
  \begin{columns}[T,onlytextwidth]
    \begin{column}{0.56\textwidth}
      \begin{lstlisting}[language=C,numbers=left,
        caption={Fungsi \texttt{run\_dfa} untuk simulasi DFA}]
int run_dfa(const DFA *d, const int *input, int n) {
    int state = d->start;
    for (int i = 0; i < n; i++) {
        state = d->delta[state][input[i]];
        if (state < 0) return 0;
    }
    return d->accept[state];
}
      \end{lstlisting}
    \end{column}
    \begin{column}{0.40\textwidth}
      \begin{itemize}
        \item \texttt{DFA} menyimpan \texttt{delta[state][simbol]},
              \texttt{start}, dan \texttt{accept}.
        \item Tabel transisi dua dimensi membuat simulasi \(O(n)\).
        \item State negatif menandakan transisi tidak terdefinisi.
        \item Simulator NFA memakai himpunan state dan
              \textit{epsilon closure}.
      \end{itemize}
    \end{column}
  \end{columns}
\end{frame}

% ---------------------------------------------------------------------
\begin{frame}{NFA: Definisi}
  \begin{block}{\textit{Nondeterministic Finite Automaton} (NFA)}
    \small
    \(M = (Q, \Sigma, \delta, q_0, F)\), dengan
    \(\delta : Q \times (\Sigma \cup \{\varepsilon\}) \to
    \mathcal{P}(Q)\). Untuk setiap state dan simbol, NFA dapat memiliki
    nol, satu, atau beberapa state tujuan.
  \end{block}

  \begin{itemize}
    \item Mesin dapat ``menebak'' langkah yang benar.
    \item String diterima jika ada setidaknya satu jalur dari \(q_0\)
          ke state penerima.
    \item NFA tidak lebih kuat daripada DFA, tetapi lebih mudah
          disusun dari ekspresi reguler.
    \item Simulasi menyimpan himpunan state, bukan satu state.
  \end{itemize}
\end{frame}

% ---------------------------------------------------------------------
\begin{frame}{Contoh NFA: String Berakhir dengan \texttt{abb}}
  \begin{center}
    \begin{tikzpicture}[
        state/.style={circle,draw=ubbblue,fill=ubbblue!7,
          minimum size=0.80cm,inner sep=0pt,font=\small},
        accepting/.style={state,double,double distance=1.4pt},
        ar/.style={-{Stealth[length=1.8mm]},draw=ubbblue!70,thick}
      ]
      \node[state] (q0) at (0,0) {\(0\)};
      \node[state] (q1) at (2.7,0) {\(1\)};
      \node[state] (q2) at (5.4,0) {\(2\)};
      \node[accepting] (q3) at (8.1,0) {\(3\)};
      \draw[ar] ([xshift=-0.75cm]q0.west) -- (q0.west);
      \draw[ar] (q0) edge[loop above] node{\scriptsize \(a,b\)} (q0);
      \draw[ar] (q0) -- node[above]{\scriptsize \(a\)} (q1);
      \draw[ar] (q1) -- node[above]{\scriptsize \(b\)} (q2);
      \draw[ar] (q2) -- node[above]{\scriptsize \(b\)} (q3);
    \end{tikzpicture}
  \end{center}

  \vspace{2pt}
  \begin{itemize}
    \item NFA menerima \((a \mid b)^{*}abb\).
    \item Di state 0, mesin dapat tetap di 0 atau pindah ke 1 saat
          membaca \(a\).
    \item Pilihan inilah inti nondeterminisme.
  \end{itemize}
\end{frame}

% ---------------------------------------------------------------------
\begin{frame}{\textit{Epsilon Closure}}
  \begin{block}{Definisi}
    \small
    \(\varepsilon\text{-closure}(S)\) adalah himpunan state yang dapat
    dicapai dari \(S\) melalui nol atau lebih transisi
    \(\varepsilon\).
  \end{block}

  \begin{itemize}
    \item Dihitung dengan penelusuran graf; kompleksitas linear terhadap
          jumlah transisi.
    \item Transisi simbol harus mempertimbangkan seluruh state dalam
          closure.
  \end{itemize}

  \vspace{3pt}
  \begin{exampleblock}{Contoh}
    \footnotesize
    Transisi \(q_0 \xrightarrow{\varepsilon} q_1\),
    \(q_1 \xrightarrow{\varepsilon} q_2\),
    \(q_2 \xrightarrow{a} q_3\).\\
    \(\varepsilon\text{-closure}(\{q_0\}) = \{q_0,q_1,q_2\}\);
    membaca \(a\) mencapai \(q_3\).
  \end{exampleblock}
\end{frame}

% ---------------------------------------------------------------------
\begin{frame}{Simulasi NFA}
  \begin{enumerate}
    \setlength{\itemsep}{3pt}
    \item Mulai dengan \(\varepsilon\text{-closure}(\{q_0\})\).
    \item Untuk setiap simbol \(a\), hitung himpunan state tujuan dari
          semua state saat ini, lalu ambil \textit{epsilon closure}-nya.
    \item Jika himpunan state kosong, tolak.
    \item Setelah masukan habis, terima jika himpunan state beririsan
          dengan \(F\).
  \end{enumerate}

  \vspace{5pt}
  \begin{block}{Kompleksitas}
    \footnotesize
    Simulasi NFA sederhana memerlukan \(O(n \cdot |Q|^2)\). Untuk lexer,
    NFA biasanya dikonversi menjadi DFA agar simulasi tetap \(O(n)\).
  \end{block}
\end{frame}

% ---------------------------------------------------------------------
\begin{frame}{\textit{Subset Construction}: NFA ke DFA}
  \begin{enumerate}
    \setlength{\itemsep}{2pt}
    \item State awal DFA adalah \(\varepsilon\text{-closure}(\{q_0\})\).
    \item Untuk setiap state DFA \(S\) yang belum diproses dan setiap
          simbol \(a\):
          \begin{enumerate}
            \item hitung
                  \(T=\varepsilon\text{-closure}(\bigcup_{q\in S}
                  \delta(q,a))\);
            \item jika \(T\) belum ada, tambahkan sebagai state DFA;
            \item tetapkan \(\delta_D(S,a)=T\).
          \end{enumerate}
    \item State DFA menerima jika \(S \cap F \neq \varnothing\).
  \end{enumerate}

  \vspace{4pt}
  {\footnotesize Jumlah state DFA paling banyak \(2^{|Q|}\); dalam
  praktik banyak himpunan tidak terjangkau sehingga hasilnya lebih
  kecil.}
\end{frame}

% ---------------------------------------------------------------------
\begin{frame}{Contoh Perhitungan \textit{Subset Construction}}
  \small
  \begin{itemize}
    \item State awal \(A=\{0\}\).
    \item Dari \(A\) membaca \(a\): \(\mathrm{move}(A,a)=\{0,1\}\),
          closure \(\{0,1\}\), jadi \(A \xrightarrow{a} B\).
    \item Dari \(A\) membaca \(b\): \(\mathrm{move}(A,b)=\{0\}\), jadi
          \(A \xrightarrow{b} A\).
    \item Dari \(B=\{0,1\}\) membaca \(a\): \(\{0,1\}\), jadi
          \(B \xrightarrow{a} B\).
    \item Dari \(B\) membaca \(b\): \(\{0,2\}\), jadi
          \(B \xrightarrow{b} C\).
  \end{itemize}

  \vspace{4pt}
  \begin{block}{Lanjutan}
    \footnotesize
    Dari \(C=\{0,2\}\) membaca \(b\) menghasilkan \(\{0,3\}=D\); dari
    \(D\) membaca \(b\) kembali ke \(A\). State penerima adalah \(D\)
    karena memuat state NFA 3.
  \end{block}
\end{frame}

% ---------------------------------------------------------------------
\begin{frame}{Tabel Konversi NFA ke DFA (Artefak Sub-CPMK 2.3)}
  \small
  \setlength{\tabcolsep}{5pt}
  \begin{tabular}{@{}c c c c c@{}}
    \toprule
    \textbf{State DFA} & \textbf{Himpunan NFA} &
    \textbf{\(a\)} & \textbf{\(b\)} & \textbf{Penerima} \\
    \midrule
    \(A\) & \(\{0\}\) & \(B\) & \(A\) & tidak \\
    \(B\) & \(\{0,1\}\) & \(B\) & \(C\) & tidak \\
    \(C\) & \(\{0,2\}\) & \(B\) & \(D\) & tidak \\
    \(D\) & \(\{0,3\}\) & \(B\) & \(A\) & ya \\
    \bottomrule
  \end{tabular}

  \vspace{6pt}
  \begin{block}{Cara mengisi}
    \footnotesize
    Setiap baris dihitung dari NFA asal: hitung
    \(\text{move}(S,a)\), ambil \textit{epsilon closure}-nya, lalu beri
    nama state baru bila himpunan itu belum muncul.
  \end{block}
\end{frame}

% ---------------------------------------------------------------------
\begin{frame}{DFA Hasil Konversi}
  \begin{center}
    \begin{tikzpicture}[
        state/.style={circle,draw=ubbblue,fill=ubbblue!7,
          minimum size=0.80cm,inner sep=0pt,font=\small},
        accepting/.style={state,double,double distance=1.4pt},
        ar/.style={-{Stealth[length=1.8mm]},draw=ubbblue!70,thick}
      ]
      \node[state] (A) at (0,0) {\(A\)};
      \node[state] (B) at (3.0,0) {\(B\)};
      \node[state] (C) at (6.0,0) {\(C\)};
      \node[accepting] (D) at (9.0,0) {\(D\)};
      \draw[ar] ([xshift=-0.75cm]A.west) -- (A.west);
      \draw[ar] (A) edge[loop above] node{\scriptsize \(b\)} (A);
      \draw[ar] (B) edge[loop above] node{\scriptsize \(a\)} (B);
      \draw[ar] (A) -- node[above]{\scriptsize \(a\)} (B);
      \draw[ar] (B) -- node[above]{\scriptsize \(b\)} (C);
      \draw[ar] (C) -- node[above]{\scriptsize \(b\)} (D);
      \draw[ar,bend left=22] (C) to node[below]{\scriptsize \(a\)} (B);
      \draw[ar,bend left=35] (D) to node[below]{\scriptsize \(a\)} (B);
      \draw[ar,bend left=26] (D) to node[below]{\scriptsize \(b\)} (A);
    \end{tikzpicture}
  \end{center}

  \vspace{2pt}
  {\footnotesize DFA ini menerima bahasa yang sama dengan NFA asal,
  yaitu seluruh string yang berakhir dengan \texttt{abb}.}
\end{frame}

% ---------------------------------------------------------------------
\begin{frame}{Minimalisasi DFA: Ekuivalensi State}
  \begin{block}{Ekuivalensi state}
    \small
    Dua state \(p\) dan \(q\) ekuivalen jika untuk setiap string \(w\),
    \(\hat{\delta}(p,w) \in F\) tepat ketika
    \(\hat{\delta}(q,w) \in F\). State yang tidak ekuivalen disebut
    terbedakan.
  \end{block}

  \begin{enumerate}
    \setlength{\itemsep}{2pt}
    \item Partisi awal: \(\{F, Q \setminus F\}\).
    \item Kelompokkan state dalam satu grup berdasarkan vektor tujuan
          transisinya terhadap grup saat ini.
    \item Ulangi sampai tidak ada grup yang terpecah.
    \item Gabungkan state dalam grup yang sama.
  \end{enumerate}
\end{frame}

% ---------------------------------------------------------------------
\begin{frame}{Tabel Minimalisasi DFA (Artefak Sub-CPMK 2.3)}
  \scriptsize
  \setlength{\tabcolsep}{4pt}
  \begin{tabular}{@{}p{1.1cm}p{3.3cm}p{2.8cm}p{6.0cm}@{}}
    \toprule
    \textbf{Iterasi} & \textbf{Partisi} &
    \textbf{State dipisah} & \textbf{Alasan} \\
    \midrule
    0 & \(\{A,B,C\},\{D\}\) & --- &
      State penerima dipisahkan dari bukan penerima \\
    1 & \(\{A,B\},\{C\},\{D\}\) & \(C\) &
      \(C \xrightarrow{b} D\), sedangkan \(A,B \xrightarrow{b}\) ke
      \(\{A,B,C\}\) \\
    2 & \(\{A\},\{B\},\{C\},\{D\}\) & \(A,B\) &
      \(A \xrightarrow{b} A\), sedangkan \(B \xrightarrow{b} C\) \\
    \bottomrule
  \end{tabular}

  \vspace{5pt}
  {\scriptsize\color{ubbblue!75}Setiap pemisahan disertai alasan
  berdasarkan transisi. Hasil akhir empat grup berarti DFA sudah
  minimal.}
\end{frame}

% ---------------------------------------------------------------------
\begin{frame}{Konstruksi Thompson: Ekspresi Reguler ke NFA}
  \small
  \setlength{\tabcolsep}{4pt}
  \begin{tabular}{@{}p{3.0cm}p{10.6cm}@{}}
    \toprule
    \textbf{Ekspresi} & \textbf{Konstruksi} \\
    \midrule
    \(\varnothing\) & State awal dan penerima tanpa jalur \\
    \(\varepsilon\) & Transisi epsilon dari awal ke penerima \\
    \(a\) & Transisi \(a\) dari awal ke penerima \\
    \(r \mid s\) & State awal baru bercabang ke NFA \(r\) dan \(s\) \\
    \(rs\) & State penerima \(r\) dihubungkan ke awal \(s\) \\
    \(r^{*}\) & Loop epsilon di sekitar NFA \(r\) \\
    \bottomrule
  \end{tabular}

  \vspace{5pt}
  \[
    \text{regex} \;\to\; \text{NFA Thompson} \;\to\;
    \text{DFA} \;\to\; \text{DFA minimal}
  \]
\end{frame}

% ---------------------------------------------------------------------
\begin{frame}{Menggabungkan Token dan \textit{Maximal Munch}}
  \begin{itemize}
    \item Buat NFA terpisah untuk setiap pattern token.
    \item Tambahkan state awal baru dengan transisi \(\varepsilon\) ke
          setiap NFA.
    \item State penerima setiap NFA diberi label jenis token.
    \item Jika beberapa token cocok pada panjang sama, label dengan
          prioritas tertinggi yang dipilih.
  \end{itemize}

  \vspace{4pt}
  \begin{block}{\textit{Maximal munch} pada DFA}
    \footnotesize
    Jalankan DFA selama masih ada transisi, sambil mengingat posisi
    terakhir di state penerima. Ketika transisi berikutnya tidak ada,
    lexer kembali ke posisi terakhir dan mengeluarkan token. Pendekatan
    ini memerlukan \eng{rollback} pada buffer masukan.
  \end{block}
\end{frame}

% ---------------------------------------------------------------------
\begin{frame}{Kegiatan Mahasiswa di Kelas: P1-B}
  \begin{block}{Simulator NFA/DFA dan tabel konversi}
    \footnotesize
    \begin{enumerate}
      \setlength{\itemsep}{1pt}
      \item Gambar NFA untuk setiap kategori token pada
            \texttt{spec/tokens.md}.
      \item Lakukan \textit{subset construction}; isi tabel konversi.
      \item Minimalisasi DFA; isi tabel partisi beserta alasan.
      \item Implementasikan simulator NFA dan DFA dalam C.
      \item Uji kedua simulator dengan 30 kasus uji dari P1-A.
      \item Pastikan NFA, DFA hasil konversi, dan DFA minimal menerima
            bahasa token yang sama.
    \end{enumerate}
  \end{block}

  \vspace{2pt}
  {\footnotesize\color{ubbblue!80}Bandingkan hasil dengan alat pembanding
  laboratorium; catat perbedaan perilaku, tetapi artefak utama tetap
  simulator C.}
\end{frame}

% ---------------------------------------------------------------------
\begin{frame}{Tugas dan Artefak}
  \begin{block}{Artefak Sub-CPMK 2.2 dan 2.3}
    Simulator NFA dan DFA dalam bahasa C, tabel konversi NFA ke DFA,
    tabel minimalisasi, serta bukti bahwa ketiganya menerima bahasa
    token yang sama.
  \end{block}

  \small
  \begin{itemize}
    \item Tabel konversi lengkap: state DFA, himpunan NFA, transisi
          setiap simbol, dan status penerima.
    \item Tabel minimalisasi memuat iterasi, partisi, dan alasan setiap
          pemisahan.
    \item Simulator C berjalan dan diuji pada 30 kasus uji P1-A.
    \item Keputusan terima/tolak NFA, DFA hasil konversi, dan DFA
          minimal identik pada semua kasus.
  \end{itemize}

  \vspace{3pt}
  {\footnotesize\color{ubbblue!75}P1-B menjadi bagian Praktikum 1
  (12\%) yang dilaporkan pada Pertemuan 5.}
\end{frame}

% ---------------------------------------------------------------------
\begin{frame}{Rangkuman dan Jembatan}
  \begin{itemize}
    \item DFA deterministik: satu transisi per state-simbol, mudah
          disimulasikan dalam \(O(n)\).
    \item NFA dapat memiliki banyak tujuan dan transisi
          \(\varepsilon\); string diterima jika ada satu jalur
          penerima.
    \item \textit{Epsilon closure} menjadi inti simulasi NFA dan
          \textit{subset construction}.
    \item \textit{Subset construction} mengubah NFA menjadi DFA;
          minimalisasi menggabungkan state ekuivalen.
    \item Simulator C, tabel konversi, dan tabel minimalisasi membuktikan
          bahasa token yang sama.
  \end{itemize}

  \vspace{4pt}
  \begin{exampleblock}{Pertemuan berikutnya}
    Pertemuan 5 mengimplementasikan lexer tangan dan lexer Flex.
    Automata yang dibangun hari ini menjadi mesin pengenal token pada
    Praktikum 1.
  \end{exampleblock}
\end{frame}

\end{document}
